记 \(h=N^{-1}\)，选择固定紧环带 \(c_0\le|\eta|\le C_0\)，使其包含 \(\varphi\) 的支撑。取有限的光滑分解 \(\varphi=\sum_\nu\varphi_\nu\)，从属于该环带上孔径如 \ref{ms:prop-6-3} 所示的各球冠。在环面 \(\mathbb R^n/(100\mathbb Z)^n\) 上，将 \ref{ms:prop-6-3} 用于 \(\varphi_\nu(hD)F\)。其 \(L^{p_c}\) 范数一致不超过 \(C_\nu\|F\|_{p_c}\)：其卷积核是 \(h^{-n}\check\varphi_\nu(\cdot/h)\) 的周期化，而后者的 \(L^1\) 范数至多为 \(\|\check\varphi_\nu\|_1\)。这里 \(\check\varphi_\nu\) 表示逆 Fourier 变换。环面在一个基本立方体上采用通常的 Lebesgue 体积测度，因此周期化核在频率 \((2\pi/100)\mathbb Z^n\) 处具有相同的乘子符号。对 \(\nu\) 求和得到
\[
 \|U(T)\varphi(hD)F\|_{L^{p_c}(I\times\mathbb T^n)}
 \le C[\log(2/h)]^{2/p_c}\|F\|_{L^{p_c}(\mathbb T^n)}.
\]

在 \(\mathbb R^n\) 上，卷积核为
\[
 K_h(T,X)=(2\pi h)^{-n}
   \int e^{i(X\cdot\eta-T|\eta|)/h}\varphi(\eta)\,d\eta .
\]
对 \(T\in I\) 和 \(|X|\ge3\)，相位 \(\Phi=X\cdot\eta-T|\eta|\) 满足 \(|\nabla_\eta\Phi|\ge |X|/3\)，且二阶及更高阶导数在环带上一致有界。使用 \(h\nabla_\eta\Phi\cdot\nabla_\eta/(i|\nabla_\eta\Phi|^2)\) 作分部积分。去掉 \(h\) 后，系数及其各阶导数在支撑上均为 \(O(|X|^{-1})\)。因此反复进行 \(M+n\) 次，其中 \(M\ge1\) 为整数，得到
\[
 |K_h(T,X)|\le C_M h^M(1+|X|)^{-M}
 \qquad (T\in I,\ |X|\ge3).
\]
将 \(\mathbb R^n\) 用单位立方体 \(Q\) 铺满，令 \(10Q\) 为边长十的同心立方体。对 \(X\in Q\) 和 \(Y\notin 10Q\)，有 \(|X-Y|\ge9/2\)。故每个立方体上的远处贡献都由同一个函数控制，即
\[
 C_M h^M\big((1+|\cdot|)^{-M}*|f|\big)(X).
\]
对 \(M>n\)，此函数的 \(L^{p_c}(I\times\mathbb R^n)\) 范数至多为 \(C_M h^M\|f\|_{p_c}\)，这是由 Young 不等式和 \(|I|=1\) 得出的。

对于近处贡献，令 \(f_Q=f\mathbf 1_{10Q}\)。将 \(Q\) 的中心平移到原点，再将 \(f_Q\) 以 \(100\) 为周期作周期化。周期输入的范数为 \(\|f\|_{L^{p_c}(10Q)}\)。周期输出采用 \(\sum_{k\in\mathbb Z^n}K_h(T,X+100k)\) 作为核。对 \(X\in Q\)、\(Y\in 10Q\) 和 \(k\ne0\)，自变量 \(X-Y+100k\) 的长度至少为 \(94.5|k|\)，因为对每个非零坐标 \(k_j\)，有 \(|X_j-Y_j+100k_j|\ge 94.5 |k_j|\)。由衰减界及 \(M>n\)，在 \(I\times Q\) 上周期输出和欧氏输出的逐点差值至多为 \(C_M h^M\|f_Q\|_1\)。\(10Q\) 的体积固定，因此由环面估计与 Hölder 不等式得到
\[
 \|U(T)P_Nf_Q\|_{L^{p_c}(I\times Q)}
 \le C[\log(2/h)]^{2/p_c}\|f\|_{L^{p_c}(10Q)}.
\]
取 \(p_c\) 次幂并对 \(Q\) 求和。立方体 \(10Q\) 的重叠次数有统一上界。再加上远处输出，由 \(\log(2/h)\) 与 \(\log(2+N)\) 可比即得结论。
